% paper/sections/10b_ccd_extensions.tex
%
% §10 界面適合格子の続き（CCD の非一様格子拡張・混合偏微分）
%
\subsection{非一様格子への拡張：座標変換法}
\label{sec:ccd_metric}

§~\ref{sec:CCD} で確立した等間隔格子 CCD（係数式~\eqref{eq:coef_CCD}）を非一様格子に拡張する
（コンパクト系の非一様格子拡張の一般論は~\cite{Lele1992}）．
物理空間 $x$（非一様）を計算空間 $\xi$（一様格子）に写像し，
CCD はすべて $\xi$ 上の等間隔演算子として定義した後，
式~\eqref{eq:transform_1st_correct},~\eqref{eq:transform_2nd_correct} の連鎖律で物理空間微分に変換する．

\begin{tcolorbox}[warnbox, title={$\partial J_x/\partial\xi$ の精度：全体精度への影響}]
$J_x$ を $\Ord{h^2}$ 中心差分で評価すると $\partial J_x/\partial\xi$ が $\Ord{h}$ 以下となり，
CCD の $\Ord{h^6}$ 精度が $\Ord{h}$–$\Ord{h^2}$ に劣化する．
座標配列 $x(\xi)$ を CCD に直接入力して $J_x|_i$ と $\partial J_x/\partial\xi|_i$ を同時評価することで解決する
（詳細・精度制限は §~\ref{sec:grid_gen} の \ref{box:grid_jx_accuracy} 参照）．
\end{tcolorbox}

\noindent\textbf{共通の格子幾何データ．}
非一様格子では $J_x,J_y,\partial_\xi J_x,\partial_\eta J_y,H_x,H_y$ を
格子生成時に定まる共通の幾何データとして扱い，CCD/FCCD，HFE，Ridge--Eikonal，
PPE，速度補正，曲率，CFL 評価が同じ値を参照する．
圧力 PPE だけが高次メトリクスを使い，曲率や粘性項が低次メトリクスを使う混在は，
表面張力・圧力・粘性の面釣合いを別々に崩すため禁止する．

\subsubsection{計算空間 $\xi$ 上でのフィルタ適用順序（DCCD/UCCD 共通）}
\phantomsection\label{sec:cls_on_nonuniform}% label retained for existing cross-references.

スペクトルフィルタ（DCCD：§~\ref{sec:dissipative_ccd}）や風上安定化（UCCD6：§~\ref{sec:uccd6_def}）を
非一様格子上で適用する際には，
\textbf{計算空間 $\xi$ 上でフィルタ／安定化項を実行し，事後にメトリクス $J_x = \partial\xi/\partial x$ を適用}する．
物理空間 $x$ 上での直接フィルタリングは不均等な格子幅により高周波減衰の一様性が崩れるため，
$\xi$ 空間一様格子上での操作が必須である．
したがってフィルタは計算空間微分 $(f'_\xi, f''_\xi)$ に作用させ，
メトリクス変換はフィルタ適用後に行う．
（移流項における具体的な構成は §~\ref{sec:advection_motivation} 参照．）

\subsection{2 次元への拡張：混合偏微分 \texorpdfstring{$D_{xy}^{(11)}$}{Dxy\^{}(11)} の計算}
\label{sec:ccd_mixed}

曲率 $\kappa_{lg}$ の計算式（式~\eqref{eq:curvature_2d}）には混合偏微分 $\phi_{xy} = \partial^2\phi/\partial x\partial y$ が不可欠である．
CCD 法は1次元のブロック Thomas 法として定式化されているため，混合微分は以下の\textbf{2段階適用}で計算する．

\subsubsection{基本的な考え方：逐次的な 1 次元 CCD の適用}

\begin{tcolorbox}[resultbox, title={混合微分 $\partial^2 f/\partial x\partial y$ の計算}]
\textbf{段階1：$x$ 方向 CCD}

各行（$j = \text{const}$）に対して $x$ 方向の 1 次元 CCD 系を適用し，
各格子点 $(i,j)$ における $x$ 方向一次微分 $g_{i,j} \equiv \partial f/\partial x \big|_{(i,j)}$ を計算する．
この操作を全行（$j = 0, 1, \ldots, N_y-1$）に対して実行する．
\[
  g_{i,j} = D_x^{(1)} f \big|_{(i,j)} \quad \Leftarrow \quad x \text{ 方向 CCD}
\]

\textbf{段階2：$g = \partial f/\partial x$ への $y$ 方向 CCD}

段階1で得た $g_{i,j}$ を，今度は\textbf{$y$ 方向の関数値}とみなして，
各列（$i = \text{const}$）に対して $y$ 方向の 1 次元 CCD 系を適用する．
\[
  \phi_{xy,\,i,j} = D_y^{(1)} g \big|_{(i,j)} = D_y^{(1)}\bigl(D_x^{(1)} f\bigr)\big|_{(i,j)}
\]

\textbf{結果：} 混合微分 $\partial^2 f/\partial x \partial y$ が全格子点で得られる．
\end{tcolorbox}

\subsubsection{精度の議論：\texorpdfstring{$O(h^6)$}{O(h\textasciicircum 6)} の維持}

各方向の CCD が $\Ord{h^6}$ 精度を持つとき，段階2での入力誤差増幅により一般には $\Ord{h^5}$ に劣化しうる
（段階1の打ち切り誤差 $E_g = \Ord{h^6}$ を $y$ 方向に 1 階差分すると 1 次低下するため）．
ただし，$f \in C^8$ のもとで打ち切り誤差場 $E_g \propto h^6 f^{(7)}$ の $y$ 方向変動が $\Ord{h^6}$ に収まるため $\Ord{h^6}$ が維持される
（具体的には $\partial_y E_g = \Ord{h^6 f^{(7,1)}} = \Ord{h^6}$ が $C^8$ 正則性から従う；厳密証明は付録~\ref{app:ccd_mixed_precision}）．
格子収束テストでは $N$ を倍にするたびに誤差が $1/64$ 倍（6次収束）かを確認すること．

\medskip

非一様格子での FVM--CCD メトリクス不整合（$J_n \neq J_f$）と $\mathcal{G}^\text{adj}$ による解決については §~\ref{sec:fvm_ccd_corrector} で詳述する（本稿の純 FCCD 構成ではこの不整合は構造的に生じない）．
ジャンプ補正付き面勾配を併用する場合は，
$G_{\Gamma}^{\mathrm{adj}}=G^{\mathrm{adj}}-B^{\mathrm{adj}}$ とし，
$B^{\mathrm{adj}}$ も同じ面平均位置・同じ物理面距離 $H_f$ で構成する．

以上で，界面適合非一様格子上の CCD 座標変換と混合偏微分の準備が整う．
次節以降では，この座標変換を FCCD と Ridge--Eikonal の非一様格子定式化へ接続する．
